\chapter{Konjunktiv II}

Konjunktiv II presents something as hypothetical, unreal, wished for, or
politely distanced. It is especially common in requests, advice, wishes, and
conditional sentences.

\section{The most important forms}

The forms of \textit{sein}, \textit{haben}, and the modal verbs should be
learned directly. They are generally preferred to constructions with
\textit{würde}.

\begin{center}
\small
\rowcolors{2}{referenceBlueBg}{white}
\begin{tabular}{@{}lccc@{}}
\toprule
 & sein & haben & werden \\
\midrule
ich & wäre & hätte & würde \\
du & wärst & hättest & würdest \\
er/sie/es & wäre & hätte & würde \\
wir & wären & hätten & würden \\
ihr & wärt & hättet & würdet \\
sie/Sie & wären & hätten & würden \\
\bottomrule
\end{tabular}
\end{center}

\begin{center}
\small
\rowcolors{2}{referenceGreenBg}{white}
\begin{tabularx}{\textwidth}{@{}>{\bfseries}l l X@{}}
\toprule
Infinitive & Konjunktiv II & Typical meaning \\
\midrule
können & könnte & could / would be able to \\
dürfen & dürfte & might / would be allowed to \\
müssen & müsste & would have to / should probably \\
sollen & sollte & should / ought to \\
wollen & wollte & wanted / would want to \\
mögen & möchte & would like \\
\bottomrule
\end{tabularx}
\end{center}

Many strong verbs form Konjunktiv II from the Präteritum stem, often with an
umlaut and an ending: \textit{kam -- käme}, \textit{fand -- fände},
\textit{ging -- ginge}, \textit{wusste -- wüsste}. In ordinary speech,
\textit{würde + infinitive} is common instead: \textit{ich würde kommen}.
Weak verbs have forms identical to their Präteritum, so \textit{würde} usually
makes the hypothetical meaning clearer: \textit{ich würde arbeiten}.

\section{Polite requests and suggestions}

Konjunktiv II makes a request less direct.

\begin{center}
\small
\rowcolors{2}{referenceYellowBg}{white}
\begin{tabularx}{\textwidth}{@{}>{\bfseries}l X@{}}
\toprule
Function & Example \\
\midrule
Request & \textit{Könnten Sie mir bitte helfen?} \\
Wish & \textit{Ich hätte gern einen Kaffee.} \\
Preference & \textit{Ich möchte lieber zu Hause bleiben.} \\
Advice & \textit{Du solltest früher schlafen gehen.} \\
Suggestion & \textit{Wir könnten mit dem Zug fahren.} \\
\bottomrule
\end{tabularx}
\end{center}

\section{Unreal conditions in the present or future}

Both clauses normally use Konjunktiv II\@. If the \textit{wenn}-clause comes
first, the finite verb of the main clause follows the comma immediately.

\begin{center}
\textit{Wenn ich mehr Zeit \textbf{hätte}, \textbf{würde} ich öfter lesen.}

\textit{Ich \textbf{käme} mit, wenn ich nicht arbeiten \textbf{müsste}.}
\end{center}

The \textit{wenn}-clause can sometimes be omitted when the condition is clear:
\textit{An deiner Stelle würde ich mit dem Arzt sprechen.}

\section{Unreal situations in the past}

Past Konjunktiv II uses \textit{hätte} or \textit{wäre} plus Partizip II\@. The
choice of auxiliary is the same as in the Perfekt.

\begin{center}
\small
\rowcolors{2}{referenceRedBg}{white}
\begin{tabularx}{\textwidth}{@{}>{\bfseries}l X@{}}
\toprule
Auxiliary & Example \\
\midrule
hätte + Partizip II
  & \textit{Wenn ich gelernt hätte, hätte ich die Prüfung bestanden.} \\
wäre + Partizip II
  & \textit{Wenn wir früher losgefahren wären, wären wir pünktlich angekommen.} \\
\bottomrule
\end{tabularx}
\end{center}

This construction describes a past condition or result that did not occur.
At B1, form the ordinary unreal past with
\textit{hätte/wäre + Partizip II}. A sequence such as
\textit{würde gemacht} is incomplete; more complex forms such as
\textit{würde gemacht haben} exist but are not needed for this pattern.

\section{Wishes and comparisons}

\begin{itemize}[leftmargin=*]
  \item Present wish: \textit{Wenn ich doch mehr Zeit hätte!}
  \item Past regret: \textit{Wäre ich doch früher gekommen!}
  \item Unreal comparison: \textit{Er tut so, als ob er alles wüsste.}
\end{itemize}

\begin{tcolorbox}[
  colback=referenceBlueBg,
  colframe=genderMasculine,
  title={Choosing the form},
  fonttitle=\bfseries
]
Use \textit{wäre}, \textit{hätte}, and the common modal forms directly. Use a
familiar one-word form such as \textit{käme} or \textit{ginge} when it sounds
natural; otherwise use \textit{würde + infinitive}. For an unreal past, use
\textit{hätte/wäre + Partizip II}.
\end{tcolorbox}
